% TOP_PROBLEM: {"id": 18, "title": "Smooth 4-Dimensional Poincaré Conjecture", "category_id": 7, "category": {"id": 7, "name": "topology", "display_name": "Topology", "description": "Properties preserved under continuous deformations.", "slug": "topology", "order_index": 7, "created_at": "2026-07-31T15:26:25.671Z"}, "status": "open"}
% FIRST_POSTED: null
% !TeX program = lualatex
% Compile twice: lualatex open_problems_500.tex
% All references and prose are embedded; no BibTeX database or external inputs.
\documentclass[11pt,a4paper]{article}
\usepackage[margin=24mm,headheight=14pt]{geometry}
\usepackage{fontspec}
\setmainfont{Latin Modern Roman}
\setsansfont{Latin Modern Sans}
\setmonofont{Latin Modern Mono}
\usepackage{amsmath,amssymb,mathtools}
\usepackage{unicode-math}
\setmathfont{Latin Modern Math}
\usepackage{microtype}
\usepackage{enumitem,xurl,needspace}
\usepackage[dvipsnames]{xcolor}
\usepackage{hyperref}
\hypersetup{unicode=true,colorlinks=true,linkcolor=MidnightBlue,urlcolor=MidnightBlue,
 pdftitle={Mathematical Research Notebook},pdfauthor={Source-annotated compilation},bookmarksnumbered=true}
\usepackage{bookmark}
\usepackage{tocloft}
\setlength{\cftsecnumwidth}{3em}
\cftsetpnumwidth{2.5em}
\cftsetrmarg{4em}
\usepackage{titlesec}
\titleformat{\section}{\Large\bfseries\raggedright}{\thesection}{0.7em}{}
\usepackage{fancyhdr}
\pagestyle{fancy}\fancyhf{}
\fancyhead[L]{\small\sffamily Open Problems Math | Research notebook}
\fancyhead[R]{\small 27 September 2026}
\fancyfoot[C]{\thepage}
\setlength{\parindent}{0pt}
\setlength{\parskip}{4pt plus 1pt}
\setlength{\emergencystretch}{3em}
\setcounter{tocdepth}{1}
\setcounter{secnumdepth}{1}
\setlist[itemize]{leftmargin=3em,itemsep=3pt,topsep=4pt,parsep=0pt}
\allowdisplaybreaks
\newcommand{\N}{\mathbb{N}}
\newcommand{\Z}{\mathbb{Z}}
\newcommand{\Q}{\mathbb{Q}}
\newcommand{\R}{\mathbb{R}}
\newcommand{\C}{\mathbb{C}}
\newcommand{\F}{\mathbb{F}}
\newcommand{\A}{\mathbb{A}}
\newcommand{\PP}{\mathbb P}
\newcommand{\E}{\mathbb{E}}
\newcommand{\eps}{\varepsilon}
\newcommand{\dd}{\,\mathrm d}
\newcommand{\sref}[2]{\hyperlink{source:#1:#2}{[S#2]}}
\newcommand{\eref}[2]{\hyperlink{extra:#1:#2}{[E#2]}}
% Turn the next switch off for a shorter reading copy. The quotations preserve
% every exactTarget field, including qualifications omitted from a short summary.
\newif\ifshowcatalogue
\showcataloguetrue
\newcommand{\cataloguescope}[1]{\ifshowcatalogue
 \paragraph{Exact scope in the supplied catalogue.}
 \begin{quote}\small #1\end{quote}\fi}
\usepackage{etoolbox}
\AtBeginEnvironment{itemize}{\small}
% Standalone entry; original section content is delimited for verification.
\begin{document}
\setcounter{section}{0}
%% BEGIN ORIGINAL SECTION CONTAINER
% ================================================================
% problemId: 18
% Canonical problem ID: 18
% exactTarget SHA-256: 6ac923176ee1a97030ac0bbf63442c011945be793d816f1e21f39f437f72c34b
\clearpage
\section*{\texorpdfstring{Smooth 4-Dimensional Poincaré Conjecture}{Smooth 4-Dimensional Poincaré Conjecture}}\label{18}
\subsection{Definitions and mathematical statement}
A closed manifold is compact without boundary. A smooth homotopy $4$-sphere is a smooth closed $4$-manifold $M$ for which there exist continuous maps $M\leftrightarrows S^4$ whose compositions are homotopic to the identities. The conjecture is
\[
 M\simeq S^4\quad\Longrightarrow\quad M\cong_{\rm diff} S^4.
\]
A diffeomorphism must be smooth with smooth inverse. A homeomorphism alone does not settle this smooth-category question.
\par\smallskip\textit{Formulation and concept references:} \sref{18}{1}.
\cataloguescope{Determine whether every smooth closed 4-manifold homotopy equivalent to the standard 4-sphere S\textasciicircum{}4 is diffeomorphic to S\textasciicircum{}4.}
\subsection{Short English statement}
Is every smooth four-dimensional space with the homotopy type of a sphere smoothly the standard sphere?
%% BEGIN RESEARCH INSERTION
\subsection{Research attempt}
\paragraph{Current frontier.}
\textit{Research check: 27 September 2026.}
Gerig's construction gives a particular ECH/Seiberg--Witten-based invariant
that is constant on homotopy four-spheres; the result is an obstruction to
that invariant, not a classification of smooth structures. His Section 1
also records the extension theorem for diffeomorphisms of $S^3$ and the
Cerf--Palais disk theorem. \sref{18}{1}.
The retrieved research does not supply a verified resolution of the full
smooth problem. The catalogue's code-review link is retained as supplied,
but no formal proof in that repository has been imported or certified in
this attempt. A claim about a formalization would require checking its
actual theorem and dependencies, not its filename.

\paragraph{Attack route.}
An invariant route needs a quantity that is both computable and sensitive
to a possible exotic sphere; a handle route needs geometric cancellation;
and a puncturing route needs control of the smooth end, not just the
homotopy type. We combine the latter two to locate exactly which data an
algebraic cancellation argument discards.

\paragraph{Attempt: the compact puncture criterion.}
Let $M$ be a smooth homotopy four-sphere and remove the interior of a
small smoothly embedded standard ball $B$. The compact manifold
$C=M\setminus\operatorname{int}B$ has boundary $S^3$ and is contractible.
Indeed, van Kampen applied with a collar to $M=C\cup B$ gives
$\pi_1(C)=\pi_1(M)=0$. In Mayer--Vietoris, the boundary map
$H_4(M,\Z)\to H_3(S^3,\Z)$ sends the oriented fundamental class to the
boundary fundamental class and is an isomorphism. Exactness, together
with the vanishing of the other reduced homology groups of $M$ and $B$,
then gives $\widetilde H_i(C,\Z)=0$ for every $i$.
A compact smooth manifold has finite CW type. If any higher homotopy group
of $C$ were nonzero, its first nonzero one would give a nonzero homology
group by Hurewicz; hence all vanish, and Whitehead's theorem gives
contractibility.

If $C$ is diffeomorphic to $D^4$, capping it with $B$ gives two standard
balls glued along an $S^3$ diffeomorphism. After adjusting orientations,
that diffeomorphism extends across a ball by the theorem quoted in
Gerig's Theorem 1.2. The double is therefore standard $S^4$.
Conversely, if $M$ is standard, Cerf--Palais (Gerig's Theorem 1.3) makes
the embedded ball ambiently isotopic to a standard small ball, so its
compact complement is a standard ball. \sref{18}{1}.
Thus the full sphere problem is equivalent to the assertion that every
smooth compact contractible four-manifold with boundary diffeomorphic to
$S^3$ is a standard ball: any such manifold can be capped to give a
homotopy sphere, by the same van Kampen and homology argument.

\paragraph{Attempt: what the handle chain complex really proves.}
Take a handle decomposition of $M$ with one 0-handle and one 4-handle;
write $a,b,c$ for the numbers of 1-, 2-, and 3-handles. The cellular
chain complex has the middle part
\[
 0\longrightarrow\Z^c\xrightarrow{\partial_3}\Z^b
       \xrightarrow{\partial_2}\Z^a\longrightarrow0.
 \tag{22.1}
\]
It is exact: $H_1=H_2=H_3=0$, while the top and bottom differentials
vanish with these oriented single extreme handles. It follows that
$b=a+c$. Since free abelian groups are projective, (22.1) splits. Choose
preimages of a basis of $\Z^a$ and a basis of $\ker\partial_2$ given by
$\partial_3$. In these bases the matrices are exactly
\[
 \partial_2=\begin{pmatrix}I_a&0\end{pmatrix},\qquad
 \partial_3=\begin{pmatrix}0\\I_c\end{pmatrix}.
 \tag{22.2}
\]
This establishes algebraic cancellation of the entire middle chain
complex with no assumption on an invariant.

It does not establish geometric handle cancellation. The matrix entries
count signed intersections. The algebraic count $+1$ is compatible with
geometric intersections of signs $+1,+1,-1$; changing bases does not by
itself provide a framed embedded Whitney disk removing the last pair.
Similarly, abelianized attaching words do not determine the fundamental
group or the attaching-link isotopy class. For a concrete warning, consider
\[
 \langle x,y\mid (xy)^2=x^3,\;x^3=y^5\rangle.
 \tag{22.3}
\]
Its exponent-sum matrix is
$\left(\begin{smallmatrix}-1&2\\3&-5\end{smallmatrix}\right)$,
with determinant $-1$, yet it has a nontrivial quotient: in $A_5$, choose
$x=(1\,2\,3)$ and $y=(1\,4\,3\,2\,5)$ (composition right to left).
Then $x^3=y^5=1$ and $xy=(1\,4)(2\,5)$, so both relators vanish.
The image is nontrivial. This example is not claimed to present the
fundamental group of a homotopy sphere. It proves only that a unimodular
exponent matrix cannot replace the nonabelian or geometric conditions.

A valid sufficient cancellation statement must instead produce actual
handle slides and isotopies for which each paired attaching sphere meets
its corresponding belt sphere in one transverse point with the required
framing, with a valid successive cancellation order. If all middle
handles cancel in this sense, only a 0- and a 4-handle remain and the
sphere is standard by the extension theorem. The algebra in (22.2)
does not produce those moves.

\paragraph{Stress test.}
Contractibility is a homotopy conclusion, not a diffeomorphism to a ball.
A diffeomorphism of the open puncture to $\R^4$ also cannot simply be
assumed to extend over a chosen compactification; Gerig explicitly
separates this end issue. \sref{18}{1}.
Nor can a high-dimensional Whitney or h-cobordism argument be applied
without checking its dimension assumptions. The algebraic cancellation
above is valid in dimension four precisely because it never asserts the
missing embedded-disk move. The finite permutation test for (22.3)
checks the advertised group-theoretic obstruction, not a Kirby diagram.

\paragraph{Outcome.}
\textbf{Outcome: Exploratory approach with unresolved gap.}
The compact-puncture reduction and the exact split chain complex are
established within their stated scope. A unimodular but nontrivial
presentation explicitly breaks the attempt to upgrade homology to
geometric cancellation. No exotic homotopy sphere and no universal
geometric cancellation procedure is supplied. A solution still needs
smooth attaching and end control that survives the loss of information
in the chain complex.
%% END RESEARCH INSERTION
\subsection{Sources}
\begin{itemize}\raggedright
\item[\textnormal{[S1]}]\hypertarget{source:18:1}{} Chris Gerig, No homotopy 4-sphere invariants using ECH = SWF, Algebraic \& Geometric Topology 21 (2021), Conjecture 1.1.
\url{https://arxiv.org/pdf/1905.10938}.
{\small\textit{Catalogue formulation source.}}
\item[\textnormal{[S2]}]\hypertarget{source:18:2}{} smooth4pc-t73-lean: Independent Review.
\url{https://github.com/toffee-desuwa/smooth4pc-t73-lean/blob/main/docs/INDEPENDENT_REVIEW.md}.
{\small\textit{Catalogue research/status reference; not a proof endorsement.}}
\item[\textnormal{[S3]}]\hypertarget{source:18:3}{} The Smooth Four-Dimensional Poincare Conjecture.
\url{https://epoch.ai/frontiermath/open-problems/smooth-poincare-4d}.
{\small\textit{Catalogue research/status reference; not a proof endorsement.}}
\end{itemize}

%% END ORIGINAL SECTION CONTAINER
\end{document}
